% 第4章 回归分析
\chapter{回归分析}

\begin{chapteroverview}
\textbf{这个分类是干什么的？}

回归分析是\textbf{建模和预测}的主力工具。它要回答的是：给定一堆自变量（客流、面积、促销天数……），因变量（销售额）怎么变？每个自变量\textbf{影响多大、方向是正是负}？能不能用这个方程去\textbf{预测}新店的销售额？相关分析只问"动不动一起变"，回归则进一步给出"变多少"的定量方程。

\textbf{美赛什么题型常用？}

回归是美赛\textbf{出镜率最高}的建模方法。C题几乎必做：用回归找关键影响因素、做预测；A/B题建立变量间的定量关系模型。论文"模型建立""结果分析"部分，回归系数表是核心证据。

\textbf{学习路径建议}
先掌握\textbf{线性回归}（一切的基础）和它的体检项\textbf{共线性VIF}；再学\textbf{曲线回归}（关系不是直线时）、\textbf{分层回归}（控制背景后看增量贡献）、\textbf{逐步回归}（变量太多时自动筛选）；因变量是\textbf{分类}时用\textbf{逻辑回归}。最后再根据数据特点认识岭回归、Lasso、Poisson、Logistic族等进阶模型。
\end{chapteroverview}

\section{基础级算法}

%% ========== 线性回归 ==========
\basicalgo{{\starmark}线性回归（最小二乘法）}{linear_reg}

\begin{algointro}
线性回归就是\textbf{找一条最佳直线（其实是多维平面）}，让它尽量穿过所有数据点，从而用自变量预测因变量。本案例用日均客流量、营业面积、月促销天数、竞争门店数预测月销售额。
\end{algointro}

\begin{algoscene}
\begin{itemize}
    \item 因变量是连续数值（销售额、温度、成绩），想找多个自变量对它的线性影响
    \item 美赛C题建立预测模型、量化各因素边际效应
    \item 论文中"影响因素分析"的核心模型
\end{itemize}
\end{algoscene}

\begin{algoprinciple}
想象散点图上一堆点，你画一条直线去"贴近"它们。怎么算最佳？让每个点到直线的\textbf{纵向距离平方和}最小——这就是"最小二乘法"。

回归方程长这样：
$$y = b_0 + b_1 x_1 + b_2 x_2 + \cdots + b_k x_k$$
\begin{itemize}
    \item $b_0$是截距（所有自变量为0时的基线值）
    \item $b_j$是$x_j$的系数：$x_j$每涨1单位，$y$平均涨$b_j$单位（控制其他变量不变）
    \item 本案例：竞争门店数系数$-2.025$，即周边每多一家竞争店，月销售额平均降2.025万元；月促销天数系数$+1.149$，每多促销1天销售额涨1.149万元
\end{itemize}

\textbf{$R^2$}是拟合优度：本案例0.788，说明四个自变量一起解释了78.8\%的销售额波动。
\end{algoprinciple}

\begin{algosposs}
\begin{enumerate}
    \item 点击菜单 \textbf{分析 → 回归 → 线性}
    \item 因变量选\textbf{月销售额}，自变量选\textbf{日均客流量、营业面积、月促销天数、竞争门店数}
    \item 方法选\textbf{输入}（Enter，全部强行进）
    \item 点击\textbf{统计量}，勾选"估计""置信区间""模型拟合度""共线性诊断""Durbin-Watson"
    \item 点击\textbf{绘制}，把\verb|*ZRESID|作$y$、\verb|*ZPRED|作$x$，画残差图查异方差
    \item 点击\textbf{确定}
\end{enumerate}

\textbf{默认参数参考}：缺失按删除（drop），含截距，协方差类型自动，VIF阈值10。本案例有效样本147行，$R^2=0.788$。
\end{algosposs}

\begin{algoresult}
\begin{itemize}
    \item \textbf{模型汇总表}：$R^2=0.788$、调整$R^2=0.782$；Durbin-Watson=1.84（接近2，残差无自相关）
    \item \textbf{ANOVA表}：$F=131.9$，$p=8.55\times10^{-47}$，模型整体显著
    \item \textbf{系数表}是核心：每个自变量一行
    \begin{itemize}
        \item \textbf{系数$B$}：方向和大小。客流$+0.02$、面积$+0.07$、促销天数$+1.15$、竞争店$-2.03$
        \item \textbf{$t$和$p$}：本案例4个变量全部$p<0.001$，都显著
        \item \textbf{95\%置信区间}：不含0即显著
    \end{itemize}
    \item \textbf{诊断前提}：本案例VIF最大仅1.03（无共线性），Breusch-Pagan $p=0.983$（无异方差），残差Shapiro $p=0.49$（正态）——前提全满足，结果可信
\end{itemize}
\end{algoresult}

\begin{algonotes}
\textbf{什么时候用？}连续因变量+多个连续/分类自变量，关系近似直线。

\textbf{什么时候不用？}
\begin{itemize}
    \item 因变量是0/1（是否逾期）：用逻辑回归
    \item 因变量是计数（次数）：用Poisson/负二项回归
    \item 散点明显弯曲：用曲线回归或加二次项
\end{itemize}

\textbf{常见坑}
\begin{itemize}
    \item \textbf{坑1}：不做前提诊断就报系数。残差异方差、自相关、共线性都会让$p$值失真
    \item \textbf{坑2}：把回归关系当因果。控制了变量也只能说"关联"，不一定是因果
    \item \textbf{坑3}：用回归方程外推到数据范围之外，结果不可靠
\end{itemize}

\textbf{和相似算法的区别}：它是回归家族的"母版"——后面所有回归都是它在不同因变量/不同数据形态下的变体。
\end{algonotes}

%% ========== 共线性 VIF ==========
\basicalgo{共线性分析（VIF诊断）}{vif_diag}

\begin{algointro}
共线性检查就是看\textbf{自变量之间是不是太像了}——如果两个自变量几乎在说同一件事，回归就分不清功劳该记给谁，系数会乱跳。本案例在做房价回归前，先查面积类变量是否互相重复。
\end{algointro}

\begin{algoscene}
\begin{itemize}
    \item 多元线性/逻辑回归\textbf{建模前}的必做体检
    \item 自变量多且可能相关（如建筑面积/套内面积/使用面积）
    \item 论文中报告VIF以证明模型稳定
\end{itemize}
\end{algoscene}

\begin{algoprinciple}
想象让两个特别像的人（套内面积和建筑面积）一起解释房价。既然它俩几乎同步变大，回归就搞不清"到底是谁把房价抬上去的"，于是把系数在两者间胡乱分配——标准误被放大、显著性忽高忽低、符号甚至反过来。

\begin{itemize}
    \item \textbf{VIF（方差膨胀因子）}：$VIF_j=\frac{1}{1-R_j^2}$，其中$R_j^2$是"用其他所有自变量预测$x_j$"得到的
    \item $VIF=1$：完全不共线；$VIF>5$：要警惕；$VIF>10$：严重共线性，必须处理
    \item 本案例建筑面积$VIF=150$、套内面积$VIF=356$——爆表！因为三者$r=1.00$几乎完全重复
\end{itemize}

\textbf{条件指数}$>30$也提示整个自变量系统病态（本案例44.1）。
\end{algoprinciple}

\begin{algosposs}
\begin{enumerate}
    \item 线性回归对话框（\textbf{分析 → 回归 → 线性}）中放好因变量和全部候选自变量
    \item 点击\textbf{统计量}，勾选\textbf{共线性诊断}
    \item 运行后查看系数表中的"容差"和"VIF"列
    \item 也可在Dabbit AI SPSS工具箱单独跑"共线性分析"，自动逐步剔除
\end{enumerate}

\textbf{默认参数参考}：预警阈值$VIF=5$，严重阈值$VIF=10$；处理方式stepwise\_remove（逐步剔除最大VIF）；缺失整行删除。
\end{algosposs}

\begin{algoresult}
\begin{itemize}
    \item \textbf{初始诊断表}：建筑面积VIF=150、套内面积=356、使用面积=252，全是"严重"
    \item \textbf{处置}：逐步剔除——先踢套内面积（VIF=356），再踢建筑面积（VIF=104）
    \item \textbf{复检}：保留变量最大VIF降到7.92（建筑层数），低于严重阈值10，共线性基本消除
    \item \textbf{决策}：三选一即可（使用面积/建筑面积/套内面积），保留业务上最易解释的那个，不要全塞进去
\end{itemize}
\end{algoresult}

\begin{algonotes}
\textbf{什么时候用？}只要做含多个自变量的回归，建模前必查。

\textbf{什么时候不用？}简单一元回归（只有一个自变量）不存在共线性问题。

\textbf{常见坑}
\begin{itemize}
    \item \textbf{坑1}：VIF高了就删重要变量。优先删\textbf{业务上重复}的，保留解释力强的；或合并成一个综合指标
    \item \textbf{坑2}：VIF高但$R^2$也高、预测效果好时，对预测用途影响不大，主要影响系数解释
\end{itemize}

\textbf{和相似算法的区别}：共线性是\textbf{自变量之间}的问题（VIF查它）；异方差、自相关是\textbf{残差}的问题（各自有专门检验），别混。
\end{algonotes}

%% ========== 曲线回归 ==========
\basicalgo{曲线回归}{curve_reg}

\begin{algointro}
散点图不是直线而是\textbf{弯的}怎么办？曲线回归就是在一堆候选曲线（对数、指数、幂函数、多项式）里挑最贴合的那条。本案例：推广预算越高，销售额涨得越慢——明显是边际递减的弯形。
\end{algointro}

\begin{algoscene}
\begin{itemize}
    \item 散点图显示非线性关系（饱和、递减、加速）
    \item 边际效用递减、S形增长、复利增长等典型曲线现象
    \item 美赛中经济/生物/增长类过程建模
\end{itemize}
\end{algoscene}

\begin{algoprinciple}
直线画不平的地方，换条曲线试试。常见几种：
\begin{itemize}
    \item \textbf{对数曲线}$y=a+b\ln x$：增长越来越慢（边际递减）——本案例正是这种
    \item \textbf{指数曲线}$y=ab^x$：按固定比例爆炸式增长
    \item \textbf{幂函数}$y=ax^b$：齐次增长
    \item \textbf{多项式（二次/三次）}：先升后降或有拐点
\end{itemize}

\textbf{怎么选？}用\textbf{AIC}做裁判：同样本上谁的AIC最小且模型简单，选谁。本案例对数曲线调整$R^2=0.991$、$AIC=141.5$，比直线（$R^2=0.83$、$AIC=584.7$）好一大截，方程$y=17.83+16.11\ln x$。
\end{algoprinciple}

\begin{algosposs}
\begin{enumerate}
    \item 点击菜单 \textbf{分析 → 回归 → 曲线估计}
    \item 因变量选\textbf{销售额}，自变量选\textbf{推广投入}
    \item 勾选候选模型：线性、对数、二次、三次、指数、幂（"全部候选"）
    \item 勾选\textbf{显示ANOVA表格}和"绘制模型图"
    \item 点击\textbf{确定}
\end{enumerate}

\textbf{默认参数参考}：模型选择准则AIC，多项式最高3次，残差诊断开启，置信水平95\%。注意对数/指数/幂函数要求$x>0$、$y>0$。
\end{algosposs}

\begin{algoresult}
\begin{itemize}
    \item \textbf{模型比较表}：每种曲线的$R^2$、$AIC$。本案例对数曲线$R^2=0.991$、$\Delta AIC=-443$（相对直线大幅为负），最优
    \item \textbf{选定方程}：$y=17.83+16.11\ln(x)$，截距和系数$p<0.001$都显著
    \item \textbf{残差诊断}：残差正态$p=0.988$、游程$p=0.101$、异方差$\rho=0.052$——都正常，曲线选得合适
    \item \textbf{解释}：$x$均值处边际效应$dy/dx\approx1.20$，即投入越大、新增一块钱带来的销售增量越小
\end{itemize}
\end{algoresult}

\begin{algonotes}
\textbf{什么时候用？}散点明显弯曲且能说出具体形态时。

\textbf{什么时候不用？}
\begin{itemize}
    \item 散点近似直线：别为了高$R^2$硬上高阶多项式，会过拟合
    \item 关系是多个自变量共同造成的：单变量曲线回归不够，要做多变量非线性回归
\end{itemize}

\textbf{常见坑}
\begin{itemize}
    \item \textbf{坑1}：为追求$R^2$盲目上三次、四次多项式，曲线在两端剧烈摆动（龙格现象）
    \item \textbf{坑2}：超出数据范围外推。本案例结论只在$x\in[1.2,\ 42.85]$内成立
\end{itemize}

\textbf{和相似算法的区别}：它是\textbf{单自变量}的曲线建模；多变量非线性关系用非线性回归或广义可加模型。
\end{algonotes}

%% ========== 分层回归 ==========
\basicalgo{分层回归}{hierarchical_reg}

\begin{algointro}
分层回归就是\textbf{按理论顺序分批次}往模型里加自变量，专门看"新加入的这批变量，能在已有基础上\textbf{额外多解释多少}"。本案例先放员工背景，再放工作资源，最后放团队氛围，逐层看增量。
\end{algointro}

\begin{algoscene}
\begin{itemize}
    \item 想证明"某组变量在控制了其他变量后仍有独立贡献"
    \item 心理学/组织行为学问卷分析：先控制人口学变量，再看核心变量
    \item 美赛中区分"背景因素"和"核心干预因素"谁更重要
\end{itemize}
\end{algoscene}

\begin{algoprinciple}
普通回归一次把所有变量丢进去，分不清谁是谁的功劳。分层回归像\textbf{盖楼}：
\begin{itemize}
    \item \textbf{第1层（基线）}：先放必须控制的背景变量（年龄、司龄、工作量），$R^2=0.414$
    \item \textbf{第2层}：加入工作资源（自主性、上级支持、发展机会），$R^2$涨到0.756，\textbf{增量}$\Delta R^2=0.342$，$\Delta F=66.7$，$p<0.001$——工作资源贡献巨大
    \item \textbf{第3层}：再加入心理安全感，$R^2=0.793$，$\Delta R^2=0.038$，$\Delta F=25.9$，$p<0.001$——仍有显著增量
\end{itemize}

\textbf{关键}：分层顺序是\textbf{按理论定的}，不是按数据$p$值。控制变量先进、核心变量后进，才能干净地看到核心变量的"净增量"。
\end{algoprinciple}

\begin{algosposs}
\begin{enumerate}
    \item \textbf{分析 → 回归 → 线性}，因变量选\textbf{敬业度 job\_satisfaction}
    \item \textbf{Block 1 of 1}放第1层变量：年龄age、司龄tenure、工作量workload，点"下一个"
    \item \textbf{Block 2 of 2}放第2层：工作自主性、上级支持、职业发展
    \item \textbf{Block 3 of 3}放第3层：心理安全感
    \item 点\textbf{统计量}勾"R平方变化""共线性诊断""Durbin-Watson"
    \item 点击\textbf{确定}
\end{enumerate}

\textbf{默认参数参考}：稳健标准误HC3，$\alpha=0.05$，VIF阈值10。
\end{algosposs}

\begin{algoresult}
\begin{itemize}
    \item \textbf{模型汇总表}：重点看\textbf{R平方变化$\Delta R^2$}和对应$\Delta F$的$p$值
    \begin{itemize}
        \item 第1层$R^2=0.414$（基线模型显著）
        \item 第2层$\Delta R^2=0.342$，$\Delta F=66.7$，$p<0.001$——工作资源显著增量
        \item 第3层$\Delta R^2=0.038$，$\Delta F=25.9$，$p<0.001$——心理安全感仍有额外贡献
    \end{itemize}
    \item \textbf{最终系数表}：用HC3稳健标准误。本案例7个变量全部显著，工作自主性$\beta=0.313$效应最大
    \item \textbf{前提诊断}：VIF最大4.31（无共线），Durbin-Watson=1.92（无自相关），残差正态$p=0.753$
\end{itemize}
\end{algoresult}

\begin{algonotes}
\textbf{什么时候用？}要回答"在控制了背景变量后，核心变量是否仍有独立解释力"。

\textbf{什么时候不用？}只是想找最优预测方程、没有明确理论层级时，用逐步回归更合适。

\textbf{常见坑}
\begin{itemize}
    \item \textbf{坑1}：分层顺序乱定。必须按"先外生后内生、先控制后核心"的理论逻辑
    \item \textbf{坑2}：只看最终$R^2$，不看每一层的$\Delta R^2$和$\Delta F$——那就失去了分层的意义
\end{itemize}

\textbf{和相似算法的区别}：它和逐步回归都在选变量，但\textbf{方向相反}——逐步是数据驱动自动选，分层是理论驱动手动分层。
\end{algonotes}

%% ========== 逐步回归 ==========
\basicalgo{逐步回归}{stepwise_reg}

\begin{algointro}
候选自变量太多（十几个甚至几十个），手动选不过来？逐步回归让计算机\textbf{按$p$值自动挑变量}：最有用的先进，没用的踢出去。本案例从12个经营指标里自动挑出5个预测月销售额。
\end{algointro}

\begin{algoscene}
\begin{itemize}
    \item 自变量很多、缺乏明确理论，想自动找出最重要的少数几个
    \item 美赛C题特征筛选、变量降维
    \item 探索性分析：先粗筛出关键变量，再深入建模
\end{itemize}
\end{algoscene}

\begin{algoprinciple}
像\textbf{选秀海选}：
\begin{itemize}
    \item \textbf{向前引入}：一个个试，把$p$最小（最显著）的变量拉进模型
    \item \textbf{向后剔除}：把模型里变得不显著的踢出去
    \item \textbf{逐步（both）}：边进边出，每加一个新变量都重新检查老变量要不要留
\end{itemize}

本案例用$p$值准则、both方向，从12个候选里最终留下5个：活跃会员数、商圈客流、客单价、促销费用、门店面积，调整$R^2=0.838$，整体$F=154.8$，$p=4.3\times10^{-56}$。
\end{algoprinciple}

\begin{algosposs}
\begin{enumerate}
    \item \textbf{分析 → 回归 → 线性}，因变量选\textbf{月销售额}，把\textbf{全部候选自变量}放入自变量框
    \item "方法"下拉从"输入"改为\textbf{逐步}
    \item 点\textbf{选项}：进入$F$的概率0.05，剔除$F$的概率0.10
    \item 点\textbf{统计量}勾"共线性诊断""模型拟合度"
    \item 点击\textbf{确定}
\end{enumerate}

\textbf{默认参数参考}：筛选准则pvalue，方向both，入选$p<0.05$、剔除$p>0.1$，最大入选20个。本案例5折交叉验证测试$R^2=0.814$。
\end{algosposs}

\begin{algoresult}
\begin{itemize}
    \item \textbf{输入/除去的变量表}：一步步显示哪些变量被纳入、哪些被踢。本案例纳入5次、剔除0次
    \item \textbf{最终模型系数表}：活跃会员数系数$0.563$（$p=3.6\times10^{-50}$）、商圈客流$0.804$、客单价$0.373$——这三个关联最强
    \item \textbf{整体}：$R^2=0.843$、调整$R^2=0.838$、$F$检验$p<0.001$
    \item \textbf{交叉验证}：测试$R^2=0.814$与训练$R^2$接近，说明模型稳定不过拟合
    \item 入选变量最大VIF=1.06，共线性可控
\end{itemize}
\end{algoresult}

\begin{algonotes}
\textbf{什么时候用？}变量多、理论不清，需要自动粗筛。

\textbf{什么时候不用？}
\begin{itemize}
    \item 有明确理论模型时：用分层/全变量回归，别让数据乱选
    \item 样本量小（事件数不足）时：逐步会过拟合
\end{itemize}

\textbf{常见坑}
\begin{itemize}
    \item \textbf{坑1}：把逐步选出的变量当"因果结论"。它只是统计相关，且路径依赖——换个样本选出来的变量可能完全不同
    \item \textbf{坑2}：不做交叉验证。训练$R^2$很高但新数据崩了，就是过拟合
\end{itemize}

\textbf{和相似算法的区别}：逐步回归=自动版变量筛选；共线性VIF是它的\textbf{前置安检}；Lasso是它的正则化升级版（见进阶部分）。
\end{algonotes}

%% ========== 逻辑回归 ==========
\basicalgo{{\starmark}逻辑回归（Logistic）}{logistic_reg}

\begin{algointro}
当因变量是\textbf{是/否}这种二分类（如是否逾期、是否购买、是否患病）时，线性回归就不灵了——逻辑回归应运而生。本案例预测客户未来12个月是否逾期。
\end{algointro}

\begin{algoscene}
\begin{itemize}
    \item 因变量是0/1二分类，想找影响因素并算概率
    \item 信用评分、医疗诊断、营销响应预测
    \item 美赛中"是否发生某事件"的概率建模与分类
\end{itemize}
\end{algoscene}

\begin{algoprinciple}
线性回归输出的是"销售额=多少万"，可能算出负数、也可能算出超出范围的值。但我们想知道的是\textbf{"逾期的概率"}，它必须在0到1之间。

逻辑回归用一个叫\textbf{logistic函数（S形曲线）}把任意实数压到0--1之间：
$$\ln\frac{p}{1-p}=b_0+b_1x_1+\cdots+b_kx_k$$
左边是"对数发生比"。
\begin{itemize}
    \item 系数$b$解释成\textbf{优势比$OR=e^b$}：本案例"近12月逾期次数"系数0.685，$OR=1.983$——每多一次历史逾期，未来逾期几率约翻倍
    \item $OR>1$风险上升，$OR<1$风险下降
\end{itemize}

本案例模型AUC=0.832，区分能力不错。
\end{algoprinciple}

\begin{algosposs}
\begin{enumerate}
    \item \textbf{分析 → 回归 → 二元Logistic}
    \item 因变量选\textbf{是否逾期}（取值0/1，事件类别=1）
    \item 协变量选\textbf{年龄、年收入、负债收入比、额度使用率、近12月逾期次数、性别}
    \item 分类变量点"分类"，把性别设为哑变量（首类为参照）
    \item 点\textbf{选项}勾"Hosmer-Lemeshow""分类图""95\%置信区间"
    \item 点击\textbf{确定}
\end{enumerate}

\textbf{默认参数参考}：事件类别1，参考类别0，$\alpha=0.05$，最大迭代100，普通MLE。本案例事件95例、非事件55例，EPV=15.8（达标）。
\end{algosposs}

\begin{algoresult}
\begin{itemize}
    \item \textbf{模型整体检验}：似然比$\chi^2=51.1$，$p<0.001$，模型整体显著；McFadden $R^2=0.259$
    \item \textbf{Hosmer-Lemeshow}：$\chi^2=8.95$，$p=0.346>0.05$——预测概率和实际吻合，\textbf{校准良好}
    \item \textbf{系数表}看$p$和$OR$：
    \begin{itemize}
        \item 显著变量：负债收入比$OR=1.035$、额度使用率$OR=1.048$、近12月逾期次数$OR=1.983$
        \item 不显著：年龄、年收入、性别（$p>0.05$）
    \end{itemize}
    \item \textbf{AUC=0.832}：接近0.5是瞎猜，接近1是完美区分；0.83属于良好
    \item Youden最优阈值0.74时灵敏度0.66、特异度0.89
\end{itemize}
\end{algoresult}

\begin{algonotes}
\textbf{什么时候用？}因变量是二分类（0/1），要找影响因素或算概率。

\textbf{什么时候不用？}
\begin{itemize}
    \item 因变量是连续值：用线性回归
    \item 因变量是多分类无序：用多项Logistic；有序：用有序Logistic
    \item 事件数太少（EPV$<10$）：结果不稳，需精简变量或用精确方法
\end{itemize}

\textbf{常见坑}
\begin{itemize}
    \item \textbf{坑1}：把$OR$当概率。$OR=2$是"几率翻倍"，不是"概率翻倍"
    \item \textbf{坑2}：只看系数$p$不看AUC。$p$显著不代表模型分类能力强
\end{itemize}

\textbf{和相似算法的区别}：线性回归预测连续值，逻辑回归预测\textbf{分类概率}；Probit是它的姊妹版（用正态分布而非logistic）。
\end{algonotes}

\section{进阶级算法速览}

%% ========== 分位数回归 ==========
\advancedalgo{分位数回归}{quantile_reg}

\begin{algobrief}
\textbf{一句话}：不只看"平均"，而是拟合\textbf{条件中位数（及任意分位数）}的回归——看自变量对高分位、低分位人群的不同影响。

\textbf{美赛场景区}：收入研究中，教育对高收入者和低收入者的回报是否不同；极端值敏感的政策效应分析。

\textbf{核心原理}：普通最小二乘 minimizes 残差平方和，只估条件均值；分位数回归最小化不对称加权绝对残差，可估$\tau=0.1/0.5/0.9$等任意分位。

\textbf{操作要点}：在Dabbit AI SPSS工具箱选"分位数回归"，指定因变量、自变量和目标分位数（如0.25、0.5、0.75）。

\textbf{结果关键}：比较不同分位点上系数的大小变化——若高分位系数远大于低分位，说明效应在分布高端更强。

\textbf{注意}：它对异方差、厚尾数据更稳健，能揭示均值回归看不到的"分布异质性"。
\end{algobrief}

%% ========== Deming回归 ==========
\advancedalgo{Deming回归（正交回归）}{deming_reg}

\begin{algobrief}
\textbf{一句话}：自变量和因变量\textbf{都有测量误差}时用的回归——普通回归假设$x$完全精确，这在仪器比对时不成立。

\textbf{美赛场景区}：两种测量方法/仪器的结果一致性比较；方法学比对实验。

\textbf{核心原理}：普通OLS只让$y$方向残差最小，Deming让\textbf{垂直于回归线方向}的总误差最小（正交距离），$x$和$y$都被当作有误差。需预先知道两者测量误差方差比。

\textbf{操作要点}：工具箱选"Deming回归"，输入两列测量值，设定误差方差比。

\textbf{结果关键}：斜率和截距；若斜率接近1、截距接近0，说明两种方法一致性好。

\textbf{注意}：它和Passing-Bablok常用于方法比对。当$x$无误差时，退化为普通OLS。
\end{algobrief}

%% ========== Probit回归 ==========
\advancedalgo{Probit回归}{probit_reg}

\begin{algobrief}
\textbf{一句话}：和逻辑回归目的一样（二分类因变量），只是把S形曲线从logistic换成\textbf{正态累积分布}。

\textbf{美赛场景区}：离散选择、阈值模型（如"努力超过某阈值才成功"）；经济学probit模型。

\textbf{核心原理}：把$p$通过标准正态的反函数（probit链接）线性化，$\Phi^{-1}(p)=b_0+b_1x_1+\cdots$。系数解释不如OR直观，通常报告边际效应。

\textbf{操作要点}：\textbf{分析 → 回归 → 有序/Probit}，或在工具箱选"Probit回归"，指定响应频率和观测值。

\textbf{结果关键}：系数符号和显著性；效应大小用边际概率（$x$每变1单位，成功概率变多少）。

\textbf{注意}：Probit和Logistic结果通常非常接近，选哪个主要看习惯和理论假设；逻辑回归的OR更好解释。
\end{algobrief}

%% ========== 非线性回归 ==========
\advancedalgo{非线性回归}{nonlinear_reg}

\begin{algobrief}
\textbf{一句话}：当关系是\textbf{没法通过变量变换变成直线}的曲线（如Logistic增长、Michaelis-Menten酶动力）时，直接拟合自定义非线性方程。

\textbf{美赛场景区}：S形生长曲线、饱和动力学、传染病传播模型（logistic增长）、物理经验公式。

\textbf{核心原理}：不做线性化，直接用迭代算法（如Levenberg-Marquardt）最小化残差平方和。需要\textbf{人为指定方程形式和初始参数猜测}。

\textbf{操作要点}：\textbf{分析 → 回归 → 非线性估计}，输入自定义模型表达式和参数初值。

\textbf{结果关键}：参数的估计值、标准误、$p$值；收敛是否成功；$R^2$和残差图。

\textbf{注意}：初值给得不好容易不收敛或收敛到局部最优——这是它和曲线回归最大的门槛。
\end{algobrief}

%% ========== 对数线性模型 ==========
\advancedalgo{对数线性模型}{loglinear_model}

\begin{algobrief}
\textbf{一句话}：用回归的思路分析\textbf{多个分类变量}的列联表——把期望频数取对数后建模，拆解主效应和交互效应。

\textbf{美赛场景区}：高维列联表（三个以上分类变量）的结构分析；问卷中多类别变量的关联分解。

\textbf{核心原理}：$\log(\text{期望频数})$可分解为各变量主效应+交互项。通过比较嵌套模型的似然比$\chi^2$，判断哪些交互项重要。

\textbf{操作要点}：\textbf{分析 → 对数线性 → 模型选择}，选分类变量和设计类型。

\textbf{结果关键}：模型拟合的似然比$\chi^2$（$p>0.05$说明模型拟合好）；显著的高阶交互项说明变量间复杂关联。

\textbf{注意}：它是卡方检验在多维表上的推广。二维表时它就回到卡方独立性检验。
\end{algobrief}

%% ========== HLM ==========
\advancedalgo{多层线性模型（HLM/混合效应）}{hlm}

\begin{algobrief}
\textbf{一句话}：数据天然\textbf{嵌套分层}时用它——学生嵌套在班级里、班级嵌套在学校里，不能把它们当独立样本。

\textbf{美赛场景区}：分层抽样数据、纵向重复测量、组织内部嵌套结构；组内相关问题。

\textbf{核心原理}：把回归系数拆成"固定效应"（大家通用）和"随机效应"（各组围绕它浮动的部分），同时处理层1（个体）和层2（群体）变量。

\textbf{操作要点}：在Dabbit AI SPSS工具箱选"多层线性模型"，指定分层结构（如学生$\to$班级）和各层变量。

\textbf{结果关键}：组内相关系数ICC（高说明必须用多层模型）；固定效应系数；随机效应方差占比。

\textbf{注意}：数据有明显嵌套/聚类时，用普通回归会低估标准误、夸大显著性——这是它和普通回归最大的区别。
\end{algobrief}

%% ========== Heckman ==========
\advancedalgo{Heckman两阶段模型}{heckman}

\begin{algobrief}
\textbf{一句话}：解决\textbf{样本选择性偏差}——你能观测到的数据本身不是随机样本（如只观测到"参加了工作"的人的工资）。

\textbf{美赛场景区}：工资决定模型（只有就业者才有工资）；自选择样本下的效应估计。

\textbf{核心原理}：第一阶段用Probit估计"入选概率"，算出逆米尔斯比$\lambda$；第二阶段把$\lambda$作为额外自变量放回回归，校正选择性偏差。

\textbf{操作要点}：工具箱选"Heckman两阶段"，指定选择方程（是否入选）和结果方程（被解释的连续变量）。

\textbf{结果关键}：$\lambda$的系数显著=存在选择性偏差；校正后核心变量系数与原OLS相比变化多少。

\textbf{注意}：它要求有工具变量（只影响是否入选、不直接影响结果变量），否则识别困难。
\end{algobrief}

%% ========== 稳健回归 ==========
\advancedalgo{稳健回归（RANSAC/M估计）}{robust_reg}

\begin{algobrief}
\textbf{一句话}：普通最小二乘会被\textbf{几个极端点带偏}，稳健回归给异常点"降权"，让直线不被 outliers 绑架。

\textbf{美赛场景区}：数据有明显异常值（长尾分布）；需要回归结果抗干扰。

\textbf{核心原理}：M估计用更温和的损失函数（如Huber），对大残差不那么敏感；RANSAC反复随机抽小点子集拟合，找"最多点支持"的模型。

\textbf{操作要点}：工具箱选"稳健回归（RANSAC/M估计）"，指定因变量和自变量，选择方法。

\textbf{结果关键}：与OLS系数对比——若稳健系数和OLS差很多，说明OLS被异常点影响了；同时看被标记为异常的样本。

\textbf{注意}：它不是删异常值，而是\textbf{降权}。数据干净时，OLS更有效。
\end{algobrief}

%% ========== Poisson回归 ==========
\advancedalgo{Poisson回归}{poisson_reg}

\begin{algobrief}
\textbf{一句话}：因变量是\textbf{计数数据}（发生次数，如事故数、来电数）时用的回归。

\textbf{美赛场景区}：一天内门店客诉次数、单位面积内事件数；计数类因变量的影响因素分析。

\textbf{核心原理}：假设因变量服从泊松分布，用对数链接$\ln(\mu)=b_0+b_1x_1+\cdots$。系数解释为"发生率比IRR"（$e^b$）。

\textbf{操作要点}：\textbf{分析 → 广义线性模型}，计数分布选Poisson；或工具箱选Poisson回归。

\textbf{结果关键}：系数/$IRR$的$p$值；重点检查\textbf{过度离散}（方差远大于均值）——若过度离散，改用负二项回归。

\textbf{注意}：泊松要求\textbf{均值=方差}。现实计数数据常过度离散，此时Poisson的标准误会被低估，负二项回归更合适。
\end{algobrief}

%% ========== PLS ==========
\advancedalgo{偏最小二乘回归（PLS）}{pls_reg}

\begin{algobrief}
\textbf{一句话}：自变量多、又严重共线、样本还不多时，PLS一边降维一边回归，两全其美。

\textbf{美赛场景区}：光谱数据、高维经济指标下的预测；共线性严重且样本少的回归。

\textbf{核心原理}：同时对$X$和$Y$提取潜变量（主成分），并保证潜变量间协方差最大——既考虑$X$内部结构，又紧扣$Y$。

\textbf{操作要点}：工具箱选"PLS偏最小二乘"，指定因变量和全部自变量，用交叉验证选最佳成分数。

\textbf{结果关键}：成分数选择（交叉验证$Q^2$）；回归系数；$R^2$和预测$R^2$。

\textbf{注意}：它是\textbf{PCA+回归的结合体}。自变量严重共线、样本量$<$变量数时，它比普通回归稳。
\end{algobrief}

%% ========== Tobit ==========
\advancedalgo{Tobit回归}{tobit_reg}

\begin{algobrief}
\textbf{一句话}：因变量被\textbf{截断/删失}时用——比如一堆人"支出为0"（其实是没花），普通回归会偏。

\textbf{美赛场景区}：消费支出（大量0值）、贷款额度、持续时间等有下界且堆积在边界的数据。

\textbf{核心原理}：假设存在一个潜变量（真实支出），观测值在0处被"左删失"。用极大似然同时估计回归和删失概率。

\textbf{操作要点}：工具箱选"Tobit回归"，指定因变量、删失方向（左删失于0）和自变量。

\textbf{结果关键}：系数符号方向；删失比例；与普通OLS对比看校正幅度。

\textbf{注意}：因变量在0处堆积很多时才用；若只是大量0但不是删失概念，考虑零膨胀模型。
\end{algobrief}

%% ========== 岭回归 ==========
\advancedalgo{岭回归（Ridge）}{ridge_reg}

\begin{algobrief}
\textbf{一句话}：共线性严重时，给系数估计"加一点惩罚"，牺牲一点点无偏性换稳定——岭回归。

\textbf{美赛场景区}：自变量高度相关、VIF爆表，又不想删变量时的回归稳定化。

\textbf{核心原理}：在最小二乘目标里加$\lambda\sum b_j^2$惩罚项，把系数往0收缩。$\lambda$越大收缩越强，用交叉验证选。

\textbf{操作要点}：工具箱选"岭回归"，指定因变量和自变量，交叉验证选正则化参数$\lambda$。

\textbf{结果关键}：岭迹图（各系数随$\lambda$的变化）；交叉验证误差最小的$\lambda$；收缩后更稳定的系数。

\textbf{注意}：岭回归\textbf{不做变量选择}（系数只会缩小不会变0）。要选变量用Lasso。
\end{algobrief}

%% ========== 负二项 ==========
\advancedalgo{负二项回归}{negbin_reg}

\begin{algobrief}
\textbf{一句话}：计数数据\textbf{过度离散}（方差远大于均值）时，它是Poisson回归的修正版。

\textbf{美赛场景区}：日就诊人数、事故次数、客户投诉数——这类数据往往方差$>$均值。

\textbf{核心原理}：在泊松基础上加一个离散参数，允许方差$=\mu+\alpha\mu^2$。$\alpha=0$就退化为Poisson。

\textbf{操作要点}：工具箱选"负二项回归"，指定计数因变量和自变量。

\textbf{结果关键}：离散参数$\alpha$显著$>0$=确实过度离散、用负二项对；发生率比IRR及其$p$值。

\textbf{注意}：先做Poisson的离差检验，确认过度离散了再上负二项；若0值特别多，改用零膨胀模型。
\end{algobrief}

%% ========== 零膨胀泊松 ==========
\advancedalgo{零膨胀泊松回归（ZIP）}{zip_reg}

\begin{algobrief}
\textbf{一句话}：计数数据里\textbf{0特别多}（既有"本来就不会发生"的 structural zero，又有"可能发生但这次没发生"的）时用它。

\textbf{美赛场景区}：一年就医次数（健康人大量0）、消费频次、事故数——0堆积异常。

\textbf{核心原理}：混合模型——一部分概率用Logistic决定"是不是0"，剩下的用Poisson决定"发生几次"。

\textbf{操作要点}：工具箱选"零膨胀泊松"，指定计数因变量、零膨胀模型的自变量和计数模型自变量。

\textbf{结果关键}：Vuong检验判断是否需要零膨胀；两部分各自的系数解读。

\textbf{注意}：先做ZIP vs Poisson的Vuong检验；如果不仅0多还过度离散，升级为零膨胀负二项（ZINB）。
\end{algobrief}

%% ========== GEE ==========
\advancedalgo{广义估计方程（GEE）}{gee}

\begin{algobrief}
\textbf{一句话}：处理\textbf{重复测量/聚类数据}的回归——同一对象被反复测多次，观测不独立时用它。

\textbf{美赛场景区}：面板数据、纵向追踪数据、嵌套在群体内的重复观测。

\textbf{核心原理}：不指定完整的随机效应分布，只指定"工作相关矩阵"（如可交换、AR-1），用准似然估计边际效应。

\textbf{操作要点}：工具箱选"广义估计方程"，指定聚类ID、因变量分布族、链接函数和工作相关结构。

\textbf{结果关键}：群体平均边际效应系数；比较不同工作相关结构下结果的稳健性。

\textbf{注意}：它和HLM都处理重复测量——GEE估\textbf{总体平均效应}，HLM估个体层面的随机效应。
\end{algobrief}

%% ========== GLM ==========
\advancedalgo{广义线性模型（GLM）}{glm}

\begin{algobrief}
\textbf{一句话}：一个\textbf{总框架}，把线性回归、逻辑回归、Poisson回归统一起来——换个"分布+链接函数"就能适应不同因变量。

\textbf{美赛场景区}：因变量既不是连续正态、也不是简单二分类时的统一建模入口。

\textbf{核心原理}：三个件——指数族分布（正态/二项/泊松/伽玛）+线性预测$\eta=Xb$+链接函数$g(\mu)=\eta$。选对组合就得到各种专用模型。

\textbf{操作要点}：\textbf{分析 → 广义线性模型}，在对话框里选分布族和链接函数。

\textbf{结果关键}：模型整体似然比检验；各系数$p$值；拟合优度（偏差/自由度）。

\textbf{注意}：逻辑回归是"二项分布+logit链接"，Poisson是"泊松+对数链接"——它们都是GLM的特例。
\end{algobrief}

%% ========== 零膨胀负二项 ==========
\advancedalgo{零膨胀负二项回归（ZINB）}{zinb_reg}

\begin{algobrief}
\textbf{一句话}：计数数据\textbf{既0特别多、又过度离散}时的终极版——ZIP的离散修正。

\textbf{美赛场景区}：大量0且方差远大于均值的计数（如慢病就诊次数、极端消费行为）。

\textbf{核心原理}：两部分混合——Logistic部分区分"本来不会发生"，负二项部分处理"发生了几次"，同时吸收过度离散。

\textbf{操作要点}：工具箱选"零膨胀负二项"，分别指定零膨胀模型和计数模型的自变量。

\textbf{结果关键}：Vuong检验支持ZINB优于负二项；两部分系数；离散参数。

\textbf{注意}：模型复杂、参数多，需要较大样本。先用ZIP/ZINB/负二项三者比较（AIC最小）再定稿。
\end{algobrief}

%% ========== Lasso ==========
\advancedalgo{Lasso回归}{lasso_reg}

\begin{algobrief}
\textbf{一句话}：岭回归的"会选变量"版本——惩罚不仅把系数缩小，还能\textbf{直接把不重要变量压到0}，实现自动变量筛选。

\textbf{美赛场景区}：自变量极多（基因、词袋、海量指标）的高维建模；需要稀疏模型。

\textbf{核心原理}：把惩罚项从$\lambda\sum b_j^2$（岭）换成$\lambda\sum|b_j|$（L1）。L1惩罚的几何形状让部分系数恰好变0。

\textbf{操作要点}：工具箱选"Lasso回归"，交叉验证选正则化参数$\lambda$。

\textbf{结果关键}：交叉验证误差最小的$\lambda$；非零系数对应的变量即"被选中"的关键变量；系数路径图。

\textbf{注意}：它和逐步回归都选变量，但Lasso是\textbf{正则化}方法，更稳定、不易过拟合，适合高维数据。
\end{algobrief}
